\section{Discussion}
The critical question is whether there is direct evidence that executive function changes the relative benefit of CBT, and whether that evidence concerns cognitive restructuring or a broader intervention.
An average benefit in a trial and a proposed difficulty using one component are not mutually exclusive; the review will test how directly the available studies connect these claims.
The cited clinical literature includes evidence favorable to CBT as well as limitations in certainty and design \citep{Safren2010,Lopez2018,LopezPinar2026}.
The author-source archive is preserved to make the motivation and prior argument inspectable, while the current protocol permits findings that challenge that argument.

The initial search is limited to database match counts and does not establish an eligible evidence set.
Title/abstract terminology, database coverage, overlapping reviews, outcome choice and missing moderator analyses may constrain the eventual conclusions.
Any inability to find a direct interaction test will be reported as an evidence gap, not as proof that CBT cannot help people with executive-function difficulties.
No individual treatment recommendation follows from this protocol.
